\chapter{METODE PENELITIAN}

\section{Jenis Penelitian}
Penelitian ini menggunakan pendekatan kualitatif dengan jenis penelitian deskriptif. Pendekatan kualitatif digunakan karena penelitian berupaya memahami makna, proses, konteks, dan perspektif informan dalam situasi alamiah, bukan untuk menguji hubungan antarvariabel secara statistik \parencite{flick2018designingqualitative,creswell2018qualitativeinquiry}. Pendekatan ini sesuai dengan tujuan penelitian, yaitu memahami strategi pemasaran Rumah Makan Nasi Gerilya pada platform GrabFood, mengidentifikasi faktor internal dan eksternal, serta merumuskan alternatif strategi pemasaran berdasarkan matriks SWOT.

Jenis penelitian deskriptif digunakan karena penelitian ini bertujuan menggambarkan keadaan objek secara sistematis, faktual, dan sesuai dengan konteks lapangan. Melalui penelitian deskriptif, peneliti dapat memaparkan kondisi strategi pemasaran, kekuatan, kelemahan, peluang, ancaman, dan alternatif strategi tanpa memanipulasi objek penelitian \parencite{creswell2018qualitativeinquiry,patton2015qualitativeevaluation}. Dengan demikian, penelitian ini tidak diarahkan untuk mencari pengaruh kausal, tetapi untuk menyusun pemahaman yang mendalam mengenai kondisi pemasaran usaha dan kebutuhan strategi yang relevan.

\section{Tempat dan Waktu Penelitian}
Penelitian ini dilakukan pada Rumah Makan Nasi Gerilya di Kota Medan dengan fokus utama pada aktivitas pemasaran usaha melalui platform GrabFood. Tempat penelitian dipilih secara sengaja karena Rumah Makan Nasi Gerilya merupakan usaha kuliner yang aktif memanfaatkan platform \textit{online food delivery}, memiliki interaksi pelanggan yang tinggi, dan menghadapi dinamika persaingan yang relevan untuk dianalisis dalam konteks strategi pemasaran. Dalam penelitian kualitatif, pemilihan lokasi dapat dilakukan berdasarkan prinsip \textit{information-rich case}, yaitu memilih kasus yang mampu memberikan informasi mendalam terhadap fokus penelitian \parencite{patton2015qualitativeevaluation,creswell2018qualitativeinquiry}.

Penelitian dilaksanakan pada April sampai Mei 2026. Pra-survei tahap 1 dilakukan melalui observasi platform GrabFood pada tanggal 04 April 2026 sebagaimana tercantum pada Lampiran~\ref{app:prasurvey1}. Pra-survei tahap 2 dilakukan melalui wawancara awal dengan pemilik usaha sebagaimana tercantum pada Lampiran~\ref{app:prasurvey2}. Pra-survei tahap 3 dilakukan melalui telaah dokumentasi internal yang diperoleh dari pihak Nasi Gerilya dan diringkas pada Lampiran~\ref{app:prasurvey3}. Setelah pra-survei, penelitian utama dilanjutkan melalui wawancara mendalam, observasi lanjutan, pengumpulan dokumentasi, analisis data, dan penyusunan laporan penelitian.

\section{Batasan Operasional}
Batasan operasional diperlukan agar ruang lingkup penelitian tetap terarah dan tidak melebar ke luar masalah yang diteliti. Penelitian ini dibatasi pada strategi pemasaran Rumah Makan Nasi Gerilya pada platform GrabFood dengan menggunakan analisis SWOT. Fokus penelitian diarahkan pada kondisi pemasaran, faktor internal, faktor eksternal, dan perumusan alternatif strategi. Batasan ini dipilih karena rumusan masalah penelitian berpusat pada belum stabilnya performa penjualan, retensi pelanggan, dan daya saing usaha pada platform GrabFood.

\begin{table}[!ht]
    \centering
    \caption{Batasan Operasional Penelitian}
    \label{tab:batasan-operasional}
    \begin{tabularx}{\textwidth}{p{3.2cm}Y}
        \toprule
        \textbf{Aspek} & \textbf{Batasan Operasional} \\
        \midrule
        Objek penelitian & Strategi pemasaran Rumah Makan Nasi Gerilya pada platform GrabFood \\
        Ruang lingkup & Kondisi pemasaran, faktor internal, faktor eksternal, dan alternatif strategi pemasaran berbasis SWOT \\
        Sumber data utama & Wawancara, observasi platform GrabFood, dokumentasi usaha, ulasan pelanggan, dan data pra-survei \\
        Batas penyajian data & Dokumentasi performa internal yang bersifat sensitif disajikan dalam bentuk indeks atau data derivatif, bukan nominal absolut \\
        Hasil yang diharapkan & Rumusan strategi pemasaran yang relevan untuk memperkuat daya saing, pengalaman pelanggan, dan peluang pembelian ulang \\
        \bottomrule
    \end{tabularx}
    \tablesource{Diolah peneliti (2026).}
\end{table}

\section{Definisi Operasional}
Definisi operasional digunakan untuk menjelaskan batas makna setiap fokus penelitian agar pengumpulan dan analisis data lebih terarah. Walaupun penelitian ini bersifat kualitatif, penegasan definisi operasional tetap diperlukan agar konsep strategi pemasaran, faktor internal, faktor eksternal, dan strategi SWOT dipahami secara konsisten selama proses penelitian \parencite{creswell2018qualitativeinquiry,miles2020qualitativedata}. Definisi operasional dalam penelitian ini disajikan pada Tabel~\ref{tab:definisi-operasional}.

\begin{table}[!ht]
    \centering
    \caption{Definisi Operasional Penelitian}
    \label{tab:definisi-operasional}
    \begin{tabularx}{\textwidth}{p{3cm}YY}
        \toprule
        \textbf{Fokus} & \textbf{Makna Operasional} & \textbf{Indikator yang Diamati} \\
        \midrule
        Strategi pemasaran & Cara usaha menarik, meyakinkan, melayani, dan mempertahankan pelanggan pada platform GrabFood & STP, menu unggulan, harga, promo, tampilan toko, foto produk, respons layanan, dan pemanfaatan fitur platform \\
        Faktor internal & Kondisi dari dalam usaha yang dapat menjadi kekuatan atau kelemahan & Rasa, kualitas produk, \textit{brand recognition}, konsistensi rasa, akurasi pesanan, kapasitas dapur, SDM, dan kesiapan saat \textit{peak hour} \\
        Faktor eksternal & Kondisi dari luar usaha yang dapat menjadi peluang atau ancaman & Program GrabFood, perilaku pelanggan, ongkir, sensitivitas harga, ulasan pelanggan, dan persaingan dengan \textit{merchant} lain \\
        Strategi SWOT & Alternatif strategi yang dirumuskan dari kombinasi faktor internal dan eksternal & Strategi SO, WO, ST, dan WT yang relevan untuk penguatan daya saing dan pembelian ulang \\
        \bottomrule
    \end{tabularx}
    \tablesource{Diolah peneliti (2026).}
\end{table}

\section{Skala Pengukuran}
Penelitian ini tidak menggunakan skala pengukuran kuantitatif seperti skala Likert, interval, atau rasio karena pendekatan yang digunakan adalah kualitatif deskriptif. Data penelitian tidak diolah untuk menghasilkan skor statistik, melainkan diklasifikasikan ke dalam kategori tematik yang sesuai dengan fokus penelitian. Kategori tersebut digunakan untuk membantu peneliti membaca pola data, membandingkan informasi antarsumber, dan menyusun faktor SWOT secara sistematis.

Skala yang digunakan dalam penelitian ini bersifat kategorikal, yaitu berupa pengelompokan data ke dalam tema strategi pemasaran, faktor internal, faktor eksternal, serta alternatif strategi SO, WO, ST, dan WT. Sebagai contoh, informasi mengenai rasa otentik dan produk unggulan dikelompokkan sebagai potensi kekuatan, sedangkan informasi mengenai konsistensi rasa, pesanan tidak sesuai, dan kendala \textit{sold out} dikelompokkan sebagai potensi kelemahan. Dengan cara ini, pengukuran dalam penelitian dipahami sebagai proses kategorisasi kualitatif, bukan sebagai pemberian nilai angka.

\section{Populasi dan Sampel Penelitian}
Dalam penelitian kualitatif, istilah populasi dan sampel tidak dipahami sebagai jumlah responden yang harus mewakili populasi secara statistik, tetapi sebagai sumber informasi yang dipilih karena relevan dengan fenomena yang diteliti \parencite{patton2015qualitativeevaluation,creswell2018qualitativeinquiry}. Populasi dalam penelitian ini adalah pihak-pihak yang berkaitan dengan aktivitas pemasaran Rumah Makan Nasi Gerilya pada platform GrabFood, sedangkan sampel penelitian berupa informan yang dipilih secara \textit{purposive sampling}. Teknik ini digunakan karena penelitian membutuhkan informan yang memahami strategi pemasaran, operasional usaha, dan pengalaman pembelian melalui GrabFood.

Informan penelitian terdiri atas informan kunci, informan pendukung, dan informan tambahan. Informan kunci adalah pemilik atau pengelola utama usaha karena memiliki informasi mengenai arah strategi, kondisi internal, dan hubungan usaha dengan platform. Informan pendukung dapat berupa admin toko atau karyawan yang terlibat dalam pemrosesan pesanan. Informan tambahan dapat berasal dari pelanggan yang pernah membeli melalui GrabFood. Kriteria informan disajikan pada Tabel~\ref{tab:kriteria-informan}.

\begin{table}[!ht]
    \centering
    \caption{Kriteria Informan Penelitian}
    \label{tab:kriteria-informan}
    \begin{tabularx}{\textwidth}{p{3cm}YY}
        \toprule
        \textbf{Jenis Informan} & \textbf{Kriteria} & \textbf{Informasi yang Diharapkan} \\
        \midrule
        Informan kunci & Pemilik atau pengelola utama yang memahami arah usaha, keputusan pemasaran, kerja sama dengan platform, dan perkembangan usaha & Strategi pemasaran, kekuatan, kelemahan, peluang, ancaman, dan arah pengembangan \\
        Informan pendukung & Admin toko, karyawan dapur, atau pihak yang terlibat langsung dalam pemrosesan pesanan dan pelayanan & Masalah operasional, akurasi pesanan, konsistensi rasa, pengelolaan \textit{sold out}, dan kondisi \textit{peak hour} \\
        Informan tambahan & Pelanggan yang pernah membeli melalui GrabFood, baik pelanggan baru maupun pelanggan yang pernah membeli ulang & Penilaian terhadap harga, promo, rasa, layanan, kepuasan, dan alasan membeli ulang atau tidak \\
        \bottomrule
    \end{tabularx}
    \tablesource{Diolah peneliti (2026).}
\end{table}

\section{Jenis Data}
Data yang digunakan dalam penelitian ini terdiri atas data primer dan data sekunder. Data primer adalah data yang diperoleh langsung dari sumber pertama di lapangan, sedangkan data sekunder adalah data pendukung yang diperoleh dari dokumen, arsip, tampilan platform, ulasan pelanggan, literatur, dan sumber tertulis lain yang relevan \parencite{flick2018designingqualitative,patton2015qualitativeevaluation}. Pembagian jenis data diperlukan agar peneliti dapat membedakan data yang berasal dari pengalaman langsung informan dengan data pendukung yang memperkuat konteks penelitian.

\begin{table}[!ht]
    \centering
    \caption{Jenis Data Penelitian}
    \label{tab:jenis-data}
    \begin{tabularx}{\textwidth}{p{2.8cm}YY}
        \toprule
        \textbf{Jenis Data} & \textbf{Sumber} & \textbf{Bentuk Data} \\
        \midrule
        Data primer & Pemilik usaha, admin atau karyawan, pelanggan, dan observasi langsung pada platform GrabFood & Hasil wawancara, catatan observasi, pengalaman layanan, dan informasi mengenai strategi pemasaran \\
        Data sekunder & Dokumen usaha, tampilan toko GrabFood, ulasan pelanggan, data pra-survei, dan referensi ilmiah & Rating, jumlah ulasan, foto menu, variasi harga, daftar promo, data indeks, tabel pra-survei, buku, jurnal, dan skripsi terdahulu \\
        \bottomrule
    \end{tabularx}
    \tablesource{Diolah peneliti (2026).}
\end{table}

\section{Teknik/Metode Pengumpulan Data}
Teknik pengumpulan data digunakan untuk memperoleh gambaran yang utuh mengenai kondisi pemasaran Rumah Makan Nasi Gerilya pada platform GrabFood. Penelitian kualitatif umumnya menggunakan lebih dari satu teknik pengumpulan data agar peneliti dapat memahami fenomena dari sisi tindakan, pengalaman, dan dokumen pendukung secara bersamaan \parencite{creswell2018qualitativeinquiry,patton2015qualitativeevaluation}. Teknik pengumpulan data dalam penelitian ini meliputi wawancara, observasi, dokumentasi, dan studi pustaka.

\begin{table}[!ht]
    \centering
    \caption{Teknik/Metode Pengumpulan Data}
    \label{tab:teknik-pengumpulan-data}
    \begin{tabularx}{\textwidth}{p{3cm}YY}
        \toprule
        \textbf{Teknik} & \textbf{Pelaksanaan} & \textbf{Data yang Dikumpulkan} \\
        \midrule
        Wawancara & Dilakukan secara mendalam kepada informan yang dipilih secara purposif & Strategi pemasaran, alasan penggunaan GrabFood, kekuatan, kelemahan, peluang, ancaman, dan rencana pengembangan \\
        Observasi & Dilakukan dengan mengamati tampilan toko pada GrabFood dan, apabila memungkinkan, alur operasional usaha & Menu, harga, foto produk, promo, rating, ulasan, ketersediaan menu, dan perbandingan dengan \textit{merchant} pesaing \\
        Dokumentasi & Dilakukan dengan mengumpulkan bukti tertulis maupun visual yang berkaitan dengan objek penelitian & Tangkapan layar, data ulasan, daftar menu, promo, data indeks, dan ringkasan pra-survei \\
        Studi pustaka & Dilakukan dengan menelaah buku, jurnal, skripsi, dan referensi ilmiah yang relevan & Landasan teori pemasaran, pemasaran digital, GrabFood, SWOT, dan penelitian terdahulu \\
        \bottomrule
    \end{tabularx}
    \tablesource{Diolah peneliti (2026).}
\end{table}

\section{Uji Persyaratan Instrumen}
Dalam penelitian kualitatif, peneliti merupakan instrumen utama karena berperan langsung dalam menentukan fokus, mengumpulkan data, menafsirkan informasi, dan menarik kesimpulan \parencite{flick2018designingqualitative,creswell2018qualitativeinquiry}. Instrumen bantu yang digunakan meliputi pedoman wawancara, pedoman observasi, dan lembar dokumentasi. Uji persyaratan instrumen dilakukan dengan memeriksa kesesuaian instrumen bantu terhadap rumusan masalah, tujuan penelitian, konsep teori, dan kebutuhan data SWOT.

Pedoman wawancara disusun agar proses penggalian data tetap terarah, tetapi tetap memberi ruang bagi jawaban informan yang berkembang. Pedoman observasi digunakan untuk mencatat aspek yang diamati pada platform GrabFood, sedangkan lembar dokumentasi digunakan untuk mengelola data visual dan data tertulis. Kelayakan instrumen diperiksa melalui kesesuaian antara aspek yang ditanyakan dengan fokus penelitian, kejelasan bahasa, urutan pertanyaan, dan kemampuan instrumen menghasilkan data yang dapat diklasifikasikan ke dalam faktor internal, faktor eksternal, serta matriks SWOT. Kisi-kisi instrumen utama disajikan pada Tabel~\ref{tab:kisi-kisi-wawancara}.

\begin{table}[!ht]
    \centering
    \caption{Kisi-Kisi Pedoman Wawancara}
    \label{tab:kisi-kisi-wawancara}
    \begin{tabularx}{\textwidth}{p{3cm}YY}
        \toprule
        \textbf{Aspek} & \textbf{Pokok Pertanyaan} & \textbf{Informan Utama} \\
        \midrule
        Profil usaha & Riwayat usaha, alasan memilih GrabFood, posisi GrabFood dalam penjualan, dan target pasar utama & Pemilik usaha \\
        Strategi pemasaran & STP, menu unggulan, harga, promo, media sosial, dan pemanfaatan fitur platform & Pemilik usaha dan admin toko \\
        Kekuatan internal & Keunggulan rasa, produk unggulan, reputasi merek, dan dukungan platform & Pemilik usaha dan karyawan \\
        Kelemahan internal & Konsistensi rasa, kesalahan pesanan, menu habis, kapasitas operasional, dan hambatan layanan & Pemilik usaha, admin toko, dan karyawan \\
        Peluang eksternal & Program GrabFood, \textit{discount partnership}, \textit{upselling}, dan peluang ekspansi & Pemilik usaha \\
        Ancaman eksternal & Pesaing sejenis, harga pesaing, ongkir, sensitivitas pelanggan, dan perubahan perilaku pembelian & Pemilik usaha dan pelanggan \\
        Pengalaman pelanggan & Persepsi terhadap harga, promo, rasa, layanan, dan alasan membeli ulang atau tidak & Pelanggan \\
        \bottomrule
    \end{tabularx}
    \tablesource{Diolah peneliti (2026).}
\end{table}

\section{Uji Persyaratan Asumsi Data}
Penelitian ini tidak menggunakan uji asumsi statistik seperti normalitas, multikolinearitas, heteroskedastisitas, atau autokorelasi karena data yang digunakan bersifat kualitatif. Persyaratan asumsi data dalam penelitian ini diarahkan pada kelayakan dan keterpercayaan data kualitatif, yaitu apakah data yang diperoleh relevan dengan fokus penelitian, bersumber dari informan yang tepat, dapat dibandingkan antarteknik, dan memiliki bukti pendukung yang memadai.

Keterpercayaan data dijaga melalui triangulasi sumber, triangulasi teknik, dan konfirmasi informasi penting kepada informan utama. Triangulasi sumber dilakukan dengan membandingkan informasi dari pemilik usaha, admin atau karyawan, pelanggan, serta data yang muncul pada platform GrabFood. Triangulasi teknik dilakukan dengan membandingkan hasil wawancara, observasi, dan dokumentasi. Konfirmasi kepada informan utama dilakukan untuk memastikan bahwa ringkasan temuan mengenai kondisi usaha, kendala operasional, dan arah strategi tidak menyimpang dari maksud informan \parencite{flick2018designingqualitative,miles2020qualitativedata,patton2015qualitativeevaluation}.

\section{Metode Analisis Data}
Metode analisis data dalam penelitian ini dilakukan secara bertahap sejak proses pengumpulan data dimulai sampai penelitian selesai. Dalam penelitian kualitatif, analisis data berlangsung secara simultan dengan pengumpulan data melalui proses menyeleksi, menata, menafsirkan, dan menarik makna dari data yang diperoleh \parencite{miles2020qualitativedata,creswell2018qualitativeinquiry}. Analisis diarahkan untuk menghasilkan gambaran mengenai kondisi pemasaran Rumah Makan Nasi Gerilya pada platform GrabFood dan untuk merumuskan strategi pemasaran melalui matriks SWOT.

Langkah analisis mengikuti alur reduksi data, penyajian data, penarikan kesimpulan, dan klasifikasi faktor SWOT. Reduksi data dilakukan dengan memilih informasi yang paling relevan dengan fokus penelitian. Penyajian data dilakukan dalam bentuk uraian naratif, tabel ringkasan, atau pengelompokan tema. Selanjutnya, temuan diklasifikasikan menjadi faktor internal dan eksternal, lalu dipetakan ke dalam kekuatan, kelemahan, peluang, dan ancaman. Faktor-faktor tersebut kemudian dikombinasikan dalam matriks SWOT untuk menghasilkan strategi SO, WO, ST, dan WT \parencite{david2023strategicmanagement,gurel2017swot}. Tahapan analisis disajikan pada Tabel~\ref{tab:tahapan-analisis}.

\begin{table}[!ht]
    \centering
    \caption{Tahapan Analisis Data Penelitian}
    \label{tab:tahapan-analisis}
    \begin{tabularx}{\textwidth}{p{2.8cm}YY}
        \toprule
        \textbf{Tahap} & \textbf{Kegiatan} & \textbf{Hasil yang Diharapkan} \\
        \midrule
        Reduksi data & Menyeleksi dan menyederhanakan data hasil wawancara, observasi, dan dokumentasi sesuai fokus penelitian & Data yang relevan dengan strategi pemasaran, faktor internal, faktor eksternal, dan isu pembelian ulang \\
        Penyajian data & Menyusun data ke dalam uraian deskriptif, tabel, atau pengelompokan tematik & Gambaran terstruktur mengenai kondisi usaha, posisi terhadap pesaing, dan masalah utama pada GrabFood \\
        Identifikasi faktor & Mengelompokkan temuan ke dalam kekuatan, kelemahan, peluang, dan ancaman & Daftar faktor SWOT yang berbasis temuan lapangan dan pra-survei \\
        Penyusunan matriks SWOT & Mengombinasikan faktor SWOT untuk menghasilkan strategi SO, WO, ST, dan WT & Alternatif strategi pemasaran untuk memperkuat daya saing dan pembelian ulang \\
        Penarikan kesimpulan & Menafsirkan hasil analisis dan merumuskan jawaban atas rumusan masalah penelitian & Kesimpulan penelitian dan saran strategis \\
        \bottomrule
    \end{tabularx}
    \tablesource{Diolah peneliti (2026).}
\end{table}

\section{Pengujian Hipotesis}
Penelitian ini tidak merumuskan hipotesis statistik karena menggunakan pendekatan kualitatif deskriptif. Dalam penelitian kualitatif, tujuan utama penelitian bukan menguji dugaan sementara melalui perhitungan statistik, melainkan memahami fenomena, menafsirkan temuan, dan merumuskan jawaban atas rumusan masalah berdasarkan data lapangan \parencite{creswell2018qualitativeinquiry,miles2020qualitativedata}. Oleh karena itu, bagian pengujian hipotesis dalam penelitian ini dipahami sebagai penegasan bahwa tidak ada hipotesis statistik yang diuji.

Sebagai gantinya, penelitian ini menguji keterkaitan temuan dengan rumusan masalah melalui proses analisis kualitatif. Temuan mengenai faktor internal digunakan untuk menjawab rumusan masalah pertama, temuan mengenai faktor eksternal digunakan untuk menjawab rumusan masalah kedua, sedangkan hasil matriks SWOT digunakan untuk menjawab rumusan masalah ketiga. Dengan demikian, validitas jawaban penelitian ditentukan oleh kesesuaian antara data lapangan, klasifikasi faktor SWOT, dan rumusan strategi yang dihasilkan.

\section{Teknik Pengujian Hipotesis}
Karena penelitian ini tidak menggunakan hipotesis statistik, teknik pengujian hipotesis seperti uji t, uji F, atau uji koefisien determinasi tidak digunakan. Teknik yang digunakan untuk menguatkan hasil penelitian adalah pemeriksaan konsistensi temuan melalui triangulasi sumber, triangulasi teknik, konfirmasi kepada informan utama, dan penelusuran kembali bukti dokumentasi. Langkah ini dilakukan agar hasil analisis tidak hanya berasal dari satu sumber informasi, tetapi didukung oleh data wawancara, observasi, dokumentasi, dan pra-survei.

Teknik pengujian dalam penelitian ini juga dilakukan melalui penelusuran logika matriks SWOT. Setiap strategi SO, WO, ST, dan WT harus dapat ditelusuri kembali kepada faktor kekuatan, kelemahan, peluang, atau ancaman yang ditemukan dalam data. Dengan cara tersebut, rumusan strategi pemasaran yang dihasilkan tetap memiliki dasar empiris dan tidak terlepas dari kondisi nyata Rumah Makan Nasi Gerilya pada platform GrabFood.
